\documentclass[11pt]{article}
\usepackage{geometry}
\geometry{letterpaper, margin=1in}
\usepackage{amsmath}
\usepackage{graphicx}
\usepackage{hyperref}
\usepackage{booktabs}

\title{\vspace{-1.5cm}\textbf{Correctness-First Speculative Decoding on Apple MLX: \\Negative Results, Ablations, and Practical Guidance}}
\author{
  \textbf{Neural Gravity Research} \\
  \texttt{research@skintern.ai}
}
\date{}

\begin{document}
\maketitle

\begin{abstract}
Speculative decoding promises to accelerate autoregressive language model inference by using a small drafter to propose tokens that a larger target verifies. We implement and rigorously validate speculative decoding on consumer edge hardware---an 8GB Apple M3 MacBook using Apple MLX. Our implementation pairs a 0.5B-parameter drafter (\texttt{Qwen2.5-0.5B-Instruct}) with a 4-bit quantized 3B-parameter target (\texttt{Qwen2.5-3B-Instruct-4bit}). We achieve \textbf{token-identical correctness} against a greedy baseline across all test prompts. However, we measure a \textbf{-51.32\% throughput reduction} in the corrected $K=5$ configuration (43.01 TPS baseline vs.\ 20.94 TPS speculative). A sweep over draft lengths $K \in \{1,2,3,5,8\}$ remains correct in all cases but shows monotonic slowdown, with the best tested setting ($K=1$) still at \textbf{-40.58\%}. We analyze the architectural bottlenecks unique to Python/MLX on Unified Memory systems and discuss why a correctness-first implementation does not currently yield speedups on this hardware class. We additionally report Test-Time A* Search results on MATH-500, finding that the current self-reflection heuristic underperforms greedy decoding (4\% vs.\ 12\% Pass@1). This paper is therefore a practical negative-result report intended to help future edge-LLM researchers avoid the same failure modes.
\end{abstract}

\section{Introduction}

Speculative decoding~\cite{leviathan2023fast,chen2023accelerating} is a well-established technique for accelerating autoregressive LLM inference. A small ``drafter'' model rapidly proposes $K$ candidate tokens, and a larger ``target'' model verifies them in a single batched forward pass. When the drafter's predictions align with the target's, multiple tokens are accepted per target call, yielding throughput improvements proportional to the acceptance rate.

However, the vast majority of speculative decoding research has been conducted on datacenter-class GPUs (A100, H100) with dedicated VRAM, high-bandwidth PCIe interconnects, and optimized CUDA kernels. A critical open question is whether the technique transfers to \textbf{consumer edge hardware}---specifically Apple Silicon devices with Unified Memory Architecture (UMA) and the MLX framework. The 8GB M3 MacBook imposes a strict $\sim$5.4GB Metal allocation limit, requiring aggressive 4-bit quantization of even 3B-parameter models.

In this paper, we provide an honest empirical assessment. We implement speculative decoding from scratch on MLX, rigorously validate it produces token-identical output to greedy baseline decoding, and measure its real-world throughput. Our key finding is a \textbf{negative result}: speculative decoding is \textit{slower} than baseline on this hardware class, even after fixing correctness, due to bottlenecks we analyze in detail. We also include a small design ablation over draft length $K$ to show that increasing speculative depth worsens throughput in the current architecture.

\section{Implementation}

\subsection{Architecture}
We use \texttt{Qwen2.5-3B-Instruct-4bit} (target, $\sim$1.8GB) and \texttt{Qwen2.5-0.5B-Instruct} (drafter, $\sim$1GB), both loaded into the Apple M3 Unified Memory via MLX. Draft speculation uses $K=5$ tokens per iteration.

\subsection{Sequential Verification for Correctness}
A critical implementation detail is that batch-verifying draft tokens through the target model (processing all $K$ drafts in one forward pass) produced \textbf{floating-point divergence} from token-by-token baseline generation in MLX. This divergence, likely caused by different Metal attention kernel codepaths for batched vs.\ single-token inputs, accumulated over many iterations and eventually flipped argmax decisions on one of the benchmark prompts.

To guarantee token-identical output, we adopt \textbf{sequential verification}: the target model evaluates each draft token one at a time, matching the exact same processing path as baseline autoregressive decoding. This eliminates all numerical divergence at the cost of forgoing batched verification speedups.

\subsection{Draft Model Re-Prefill}
The drafter reconstructs its full KV cache from scratch each iteration by re-processing the entire accepted prefix. This is necessary because MLX's cache trimming alone proved insufficient to maintain draft-target synchronization across branches.

\section{Experiments}

All experiments were run natively on an Apple M3 MacBook (8GB RAM). Timing uses \texttt{time.perf\_counter}. Token counts are fixed at 64 per run. The main speculative benchmark is saved to \texttt{reports/benchmark\_uma\_cascade.json}; the draft-length ablation is saved to \texttt{reports/speculative\_k\_sweep.json}.

\subsection{Speculative Decoding Benchmark}

\begin{table}[h]
\centering
\begin{tabular}{@{}lccc@{}}
\toprule
\textbf{Method} & \textbf{Avg TPS} & \textbf{Speedup} & \textbf{Correctness} \\ \midrule
Baseline (Greedy)   & 43.01 & ---      & Reference   \\
Speculative ($K=5$) & 20.94 & $-51.32\%$ & Token-identical \\ \bottomrule
\end{tabular}
\vspace{2mm}
\caption{Measured speculative decoding performance from \texttt{reports/benchmark\_uma\_cascade.json}. Five prompts, all producing token-identical output (0 mismatches).}
\label{tab:spec}
\end{table}

\subsection{Draft-Length Ablation}

\begin{table}[h]
\centering
\begin{tabular}{@{}cccc@{}}
\toprule
\textbf{$K$} & \textbf{Baseline TPS} & \textbf{Speculative TPS} & \textbf{Speedup} \\ \midrule
1 & 42.25 & 25.10 & $-40.58\%$ \\
2 & 44.33 & 24.96 & $-43.70\%$ \\
3 & 44.61 & 24.15 & $-45.87\%$ \\
5 & 44.38 & 21.58 & $-51.37\%$ \\
8 & 44.68 & 19.14 & $-57.17\%$ \\ \bottomrule
\end{tabular}
\vspace{2mm}
\caption{Draft-length sweep from \texttt{reports/speculative\_k\_sweep.json}. Every tested setting remained token-identical to greedy decoding, but throughput worsened monotonically as $K$ increased.}
\label{tab:ksweep}
\end{table}

\subsection{Test-Time A* Search (TTA*) on MATH-500}

\begin{table}[h]
\centering
\begin{tabular}{@{}lccc@{}}
\toprule
\textbf{Method} & \textbf{Pass@1} & \textbf{Total Tokens} & \textbf{Wall Time} \\ \midrule
Greedy Decoding  & 12/100 (12\%) & 19,989 & 474.1s \\
TTA* Search      & 4/100 (4\%)   & 4,100  & 920.4s \\ \bottomrule
\end{tabular}
\vspace{2mm}
\caption{Test-Time A* Search results on MATH-500. All numbers from \texttt{final\_tta\_benchmark\_results.json}. TTA* underperforms greedy decoding on both accuracy and latency.}
\label{tab:tta}
\end{table}

\section{Analysis of Negative Results}

\subsection{Why Speculative Decoding is Slower}

The measured slowdown stems from three compounding factors:

\begin{enumerate}
    \item \textbf{Draft re-prefill cost}: The drafter re-processes the entire accepted prefix from scratch each iteration. As the prefix grows, this cost grows linearly.
    \item \textbf{Sequential verification}: To guarantee correctness, the target processes each draft token individually rather than in batch. This eliminates the core parallelism advantage of speculative decoding.
    \item \textbf{Python dispatch overhead}: Each token verification requires a Python$\to$Metal round-trip. With $K=5$ drafts per iteration, this adds 5 sequential dispatch calls versus 1 in baseline.
\end{enumerate}

Table~\ref{tab:ksweep} provides an additional practical lesson: increasing $K$ hurts performance monotonically in the current architecture. This is exactly the opposite of the usual speculative-decoding intuition on datacenter GPUs, and it indicates that per-iteration draft construction plus sequential target verification dominate any acceptance-rate benefit.

On datacenter GPUs, batched verification often works correctly and efficiently due to mature deterministic kernels. On MLX/Metal, the attention codepath appears to differ between batched and single-token modes, producing numerical divergence. Fixing this at the framework level (or introducing deterministic batch behavior) would be necessary before standard speculative decoding can recover its intended speedup on Apple Silicon.

\subsection{Why TTA* Underperforms Greedy}

TTA* achieved only 4\% Pass@1 vs.\ greedy's 12\%, using fewer total tokens (4,100 vs.\ 19,989) but requiring $2\times$ wall time (920s vs.\ 474s). The failure mode is clear: the self-reflection heuristic $h(n)$---which asks the model to rate its own reasoning---is unreliable on a 4-bit quantized 3B model. The model cannot accurately judge its own partial derivations, leading to systematic pruning of correct branches.

\section{Practical Guidance for Edge-LLM Researchers}

The primary value of these results is operational:

\begin{enumerate}
    \item \textbf{Validate correctness before timing.} A speculative implementation can appear to work on several prompts while silently diverging later.
    \item \textbf{Do not assume batch verification is numerically equivalent to sequential verification on MLX.} Check token identity directly.
    \item \textbf{Treat draft-cache reuse with suspicion on edge runtimes.} A slower but explicit re-prefill may be necessary to preserve correctness.
    \item \textbf{Measure $K$ instead of inheriting values from GPU literature.} On this stack, larger $K$ made performance strictly worse.
    \item \textbf{Separate prototype claims from artifact-backed results.} Saved JSON artifacts were essential in preventing incorrect conclusions.
\end{enumerate}

\section{Reproducibility}

The main speculative benchmark can be reproduced with:
\begin{quote}
\texttt{./venv/bin/python3 benchmark\_uma\_cascade.py}
\end{quote}

The TTA* artifact discussed in Table~\ref{tab:tta} is saved in:
\begin{quote}
\texttt{final\_tta\_benchmark\_results.json}
\end{quote}

The draft-length sweep is saved in:
\begin{quote}
\texttt{reports/speculative\_k\_sweep.json}
\end{quote}

\section{Conclusion}

We present measured results for speculative decoding and test-time search scaling on consumer Apple Silicon. Both techniques, in their current form on this hardware and software stack, produce negative results. Speculative decoding is correct but substantially slower; TTA* is both slower and less accurate than greedy decoding. The useful contribution is therefore not a new performance record, but a concrete, reproducible characterization of where correctness-first edge implementations fail. We expect these results to be most useful as a baseline and debugging guide for future MLX/Metal inference work.

\begin{thebibliography}{9}
\bibitem{leviathan2023fast} Leviathan, Y., Kalman, M., \& Matias, Y. (2023). Fast inference from transformers via speculative decoding. \textit{ICML 2023}.
\bibitem{chen2023accelerating} Chen, C., et al. (2023). Accelerating large language model decoding with speculative sampling. \textit{arXiv:2302.01318}.
\end{thebibliography}

\end{document}
